\section{Выводы}

В ходе выполнения второй части работы изучены просодические
закономерности речевого сигнала.

\begin{enumerate}
	\item Просодика описывает ударение и интонацию. Акцентная группа
	состоит из предъядра, ядра (главного ударного гласного) и заядра;
	интонационный тип конечной акцентной группы определяется положением
	мелодического максимума $F_0$ относительно ядра: сдвиг влево
	соответствует завершению (утверждению), в центр --- восклицанию,
	вправо --- незавершённости или вопросу.

	\item При перестановке слов <<Саша>>, <<кушал>>, <<кашу>> фраза
	сохраняет три акцентные группы и звучит примерно одинаково
	($(+)\ (+)\ (+)$). Фраза, произнесённая с ударением на первое слово,
	образует одну акцентную группу $(+\ -\ -)$ и звучит как утверждение;
	перенос ударного слова в конец даёт вид $(-\ -\ +)$, и фраза
	воспринимается как вопрос.

	\item Ударение проявляется комплексом акустических признаков:
	пиковый уровень ударного слова на 7--21~дБ, а средний уровень --- до
	20~дБ выше, чем у безударных; ударные слова имеют повышенную
	$F_0$ (145--167~Гц против 112--133~Гц) и, как правило, большую
	длительность. Совокупность этих признаков и создаёт эффект
	выделенности слога.

	\item Интонационные типы различаются положением и направлением
	движения мелодического максимума. При завершённости максимум
	расположен в начале фразы (0,24 длительности), тон к концу понижается
	на 4,7~полутона. При вопросе максимум смещён к концу (0,94), тон
	повышается на 6,7~полутона. Восклицание характеризуется высоким
	начальным тоном с последующим падением, незавершённость ---
	промежуточным контуром без терминального падения. Эмоциональная
	окраска изменяет прежде всего диапазон изменения $F_0$: при радости
	подъёмы и спады тона выражены сильнее, чем при огорчении.

	\item Методами наклона тона (аналог Pitch Bend), постоянного сдвига
	тона (Pitch Shift), изменения амплитуды (Fade) и длительности
	(Time Stretch) синтезированы варианты фразы <<Студент писал
	реферат>> с интонацией завершённости и вопроса и с ударением на
	каждое из трёх слов. Ударное слово во всех вариантах становится
	локальным максимумом контура $F_0$ (превышение 2--7~полутонов),
	а усиление амплитуды на 5~дБ и удлинение в 1,35 раза делают
	фразовое ударение более отчётливым. Показано, что Pitch Bend
	формирует интонацию, тогда как Pitch Shift лишь изменяет высоту
	голоса, сохраняя мелодический рисунок.

	\item Таким образом, просодические характеристики --- ударение и
	интонация --- являются важнейшими средствами передачи смысловой и
	эмоциональной информации в речевом сигнале. Полученные записи и
	результаты могут быть использованы при выполнении следующих работ,
	связанных с синтезом речи.
\end{enumerate}

\begin{thebibliography}{1}
	\bibitem{lobanov}
		Лобанов, Б.~М. Речевой интерфейс интеллектуальных систем :
		учеб. пособие / Б.~М.~Лобанов, В.~А.~Ефимов ; под ред.
		Б.~М.~Лобанова. --- Минск : БГУИР, 2006. --- 152 с.
	\bibitem{rabiner}
		Рабинер, Л.~Р. Цифровая обработка речевых сигналов /
		Л.~Р.~Рабинер, Р.~В.~Шафер ; пер. с англ. --- М. : Радио и
		связь, 1981. --- 495 с.
\end{thebibliography}
